% Complete proof from direct/PROOF.md; literal text, no Markdown math interpretation.
\section{Residue-one orientation and selectively even Pell indices}



\subsection*{0. Fixed equation and scope}



The principal polynomial is the sum of the squares of all 290 literal quadratic residuals in \texttt{integer442/\allowbreak{}equation.json}. Its variable order is the order in that file. It has 442 unrestricted integer variables, exact degree four, 4,652 collected monomials, and coefficient height 128,446. Every collected coefficient is supplied in \texttt{integer442/\allowbreak{}expanded\_\allowbreak{}polynomial.json}. The companion \texttt{integer446/\allowbreak{}} has 446 integer variables, 291 residuals, 4,775 monomials and the same height. These are ordinary fixed integer polynomials, not equations with exponentiation primitives.



The 442-variable source is \texttt{../\allowbreak{}architecture/\allowbreak{}unified\_\allowbreak{}derived\_\allowbreak{}shell/\allowbreak{}raw/\allowbreak{}}; the 446-variable source is \texttt{../\allowbreak{}architecture/\allowbreak{}unified\_\allowbreak{}slack\_\allowbreak{}shell/\allowbreak{}raw/\allowbreak{}}. The former uses a derived label range, the latter an additional local label coordinate. Their ordinary natural-domain arithmetic proof is \texttt{../\allowbreak{}architecture/\allowbreak{}UNIFIED\_\allowbreak{}PRODUCT\_\allowbreak{}PROOF.md}. It identifies the factorial's binomial with the reciprocal prime-power product:



\begin{verbatim}
t = t0+16, rho=F^2, B=binom(F^2+t,t), V0=(F^2+1)^t,
F B=V0+s, B=r+s+1, K=2R, N=B^K, q=N^K.
\end{verbatim}

Here these formulas name intended meanings; powers and binomials are not assumed before their certification. The actual equations have 18 power requests and 15 seed indices. The new compiler keeps the t and K seeds undoubled, doubles the other thirteen, and leaves the derived t+\allowbreak{}1 pair as its original one-step Pell addition. Every one of the eleven \texttt{(q+\allowbreak{}1)\textasciicircum{}\allowbreak{}k} requests is among the doubled seeds. The two other doubled seeds certify q\textasciicircum{}\allowbreak{}(y+\allowbreak{}1) and q\textasciicircum{}\allowbreak{}(x+\allowbreak{}1).



The theorem is an ordinary proof, formalizable in ZFC, that the new polynomial has an integer zero iff the inherited direct source has a natural zero. Composing with that source's numerical C222/\allowbreak{}refutation calibration gives



\begin{verbatim}
ZFC proves: [there is no z in Z^442 with P442(z)=0] iff Con(ZFC).
\end{verbatim}

This is not a proof of consistency, a decision of the equation's solvability, a globally smallest equation, or a newly emitted full Lean/\allowbreak{}object-ZFC proof. The named calibration interfaces remain KEndpoint.c222\_\allowbreak{}iff\_\allowbreak{}zf\_\allowbreak{}refutation and ConstructibleChoice.zf\_\allowbreak{}refutation\_\allowbreak{}iff\_\allowbreak{}zfc\_\allowbreak{}refutation. The arithmetic route does not simulate K215 and does not instantiate Jones's universal polynomial.



\subsection*{1. The two new interface ideas}



\subsubsection*{1.1 Residue one orients a signed square root}



Suppose q\textgreater{}\allowbreak{}=\allowbreak{}3, k\textgreater{}\allowbreak{}=\allowbreak{}1 and m is an integer satisfying



\begin{verbatim}
m^2=(q+1)^(2k),  m=1 (mod q).
\end{verbatim}

The first equality gives m=\allowbreak{}+\allowbreak{}(q+\allowbreak{}1)\textasciicircum{}\allowbreak{}k or m=\allowbreak{}-(q+\allowbreak{}1)\textasciicircum{}\allowbreak{}k. Their residues modulo q are +\allowbreak{}1 and -1. Since q does not divide 2, only the positive sign is possible. Therefore m=\allowbreak{}(q+\allowbreak{}1)\textasciicircum{}\allowbreak{}k. This needs no nonnegativity guard for m.



The first materialized variant adds the single quadratic row m-1-q\textasteriskcentered{}\allowbreak{}u=\allowbreak{}0 for each of eleven outputs, with signed u. The preferred variant builds the same residue condition into an existing remainder. The actual front equation has one of the forms



\begin{verbatim}
m=L*(C+q*h)+r,     or     m=N*L*h+r,
L=r+s+1.
\end{verbatim}

Replace the old remainder r everywhere by 1+\allowbreak{}q\textasteriskcentered{}\allowbreak{}u, reusing its coordinate as u. The equations become



\begin{verbatim}
m=L*(C+q*h)+1+q*u,  or  m=N*L*h+1+q*u,
L=q*u+s+2.
\end{verbatim}

All changed residuals remain quadratic: the old remainder occurs linearly in every source row, and q is an existing coordinate. The certificate lists every occurrence and its exact substitution. Once L is proved to be a positive power of q, the first equation gives m=\allowbreak{}1 modulo q.



This substitution is NOT a bijective change of arbitrary integer coordinates. It is an existence-preserving restriction justified at the semantic interface. For completeness, a genuine terminal power has m=\allowbreak{}(q+\allowbreak{}1)\textasciicircum{}\allowbreak{}k=\allowbreak{}1 modulo q and the old divisor L is q\textasciicircum{}\allowbreak{}j with j\textgreater{}\allowbreak{}=\allowbreak{}1. Thus its old bounded remainder satisfies r=\allowbreak{}1 modulo q, 0\textless{}\allowbreak{}=\allowbreak{}r\textless{}\allowbreak{}L, hence r\textgreater{}\allowbreak{}=\allowbreak{}1; u=\allowbreak{}(r-1)/\allowbreak{}q is a nonnegative integer. For soundness, u\textgreater{}\allowbreak{}=\allowbreak{}0 and q\textgreater{}\allowbreak{}0 give the old nonnegative remainder r=\allowbreak{}1+\allowbreak{}q\textasteriskcentered{}\allowbreak{}u. The recovery of the genuine powers below establishes exactly the required interface.



\subsubsection*{1.2 Weighted interval guards share one sign anchor}



For q\textgreater{}\allowbreak{}0 suppose several pairs satisfy



\begin{verbatim}
q*u_i+s_i=L-2.
\end{verbatim}

Retain u\_\allowbreak{}0\textgreater{}\allowbreak{}=\allowbreak{}0 and s\_\allowbreak{}0\textgreater{}\allowbreak{}=\allowbreak{}0. They imply L-2\textgreater{}\allowbreak{}=\allowbreak{}0. For any other pair,



\begin{verbatim}
[u_i>=0 and s_i>=0] iff u_i*s_i>=0.
\end{verbatim}

Indeed, if either were negative, a nonnegative product would make the other nonpositive, and q\textasteriskcentered{}\allowbreak{}u\_\allowbreak{}i+\allowbreak{}s\_\allowbreak{}i would be strictly negative. This contradicts L-2\textgreater{}\allowbreak{}=\allowbreak{}0. The reverse direction is immediate. This weighted lemma is the exact adaptation of the incoming Structural Compression shared-remainder lemma to the new remainder coordinates. Seven literal applications replace fourteen individual sign predicates by seven product predicates. Every pair, retained anchor, common-bound identity, and row is recorded. The sign of q is established before using this lemma; it does not depend on a binomial interpretation.



All sign predicates G\textgreater{}\allowbreak{}=\allowbreak{}0 are represented by one quadratic row



\begin{verbatim}
G=a^2+b^2+c(c+1),  a,b,c in Z.
\end{verbatim}

This expression takes only nonnegative integer values. Conversely 4G+\allowbreak{}1 is not of the excluded three-square shape; Legendre's three-square theorem and parity put its representation in this form. G is affine or quadratic here, so the final residuals remain quadratic. There are 38 such predicates in the principal source and 39 in the companion, including a strict common-ordinate predicate.



\subsection*{2. Pell notation and facts used in the compiler}



For a\textgreater{}\allowbreak{}=\allowbreak{}2 set D=\allowbreak{}a\textasciicircum{}\allowbreak{}2-1 and define



\begin{verbatim}
(a+sqrt(D))^j=chi_a(j)+psi_a(j)*sqrt(D),  j>=0.
\end{verbatim}

They are integer sequences satisfying the usual multiplication recurrence, with (chi(0),psi(0))=\allowbreak{}(1,0) and (chi(1),psi(1))=\allowbreak{}(a,1). Both are strictly increasing after the zero endpoint, and psi\_\allowbreak{}a(j)\textgreater{}\allowbreak{}=\allowbreak{}j. Every integral solution of x\textasciicircum{}\allowbreak{}2-D\textasteriskcentered{}\allowbreak{}y\textasciicircum{}\allowbreak{}2=\allowbreak{}1 has |x|=\allowbreak{}chi\_\allowbreak{}a(i), |y|=\allowbreak{}psi\_\allowbreak{}a(i) for one i\textgreater{}\allowbreak{}=\allowbreak{}0. To prove this, reduce a positive norm-one unit by powers of a+\allowbreak{}sqrt(D); the smallest nontrivial positive integral ordinate is 1 and its abscissa is a. The standard reduction therefore leaves only the identity. Independent signs of the coordinates are allowed.



Two elementary facts supply the finite indexed-Pell relation.



\begin{enumerate}

\item If i,h\textgreater{}\allowbreak{}=\allowbreak{}1 and psi\_\allowbreak{}a(i)\textasciicircum{}\allowbreak{}2 divides psi\_\allowbreak{}a(h), then i divides h and psi\_\allowbreak{}a(i) divides h. Write h=\allowbreak{}li+\allowbreak{}r with 0\textless{}\allowbreak{}=\allowbreak{}r\textless{}\allowbreak{}i, x=\allowbreak{}chi\_\allowbreak{}a(i), y=\allowbreak{}psi\_\allowbreak{}a(i). Modulo y the addition law gives psi\_\allowbreak{}a(h)=\allowbreak{}x\textasciicircum{}\allowbreak{}l psi\_\allowbreak{}a(r). The norm gives gcd(x,y)=\allowbreak{}1; monotonicity forces r=\allowbreak{}0. Expansion modulo y\textasciicircum{}\allowbreak{}2 then gives psi\_\allowbreak{}a(li)=\allowbreak{}l\textasteriskcentered{}\allowbreak{}x\textasciicircum{}\allowbreak{}(l-1)\textasteriskcentered{}\allowbreak{}y, so y divides l, hence y divides h. Also i\textless{}\allowbreak{}=\allowbreak{}h.

\item For 1\textless{}\allowbreak{}=\allowbreak{}i\textless{}\allowbreak{}=\allowbreak{}h, if chi\_\allowbreak{}a(j)=\allowbreak{}+\allowbreak{}/\allowbreak{}-chi\_\allowbreak{}a(i) modulo chi\_\allowbreak{}a(h), then j=\allowbreak{}+\allowbreak{}/\allowbreak{}-i modulo 2h. Modulo u=\allowbreak{}chi\_\allowbreak{}a(h), the norm-unit doubling identity gives (a+\allowbreak{}sqrt(D))\textasciicircum{}\allowbreak{}(2h)=\allowbreak{}-1. Reduce j modulo 2h and reflect its remainder about h. This leaves +\allowbreak{}/\allowbreak{}-chi\_\allowbreak{}a(r) with 0\textless{}\allowbreak{}=\allowbreak{}r\textless{}\allowbreak{}=\allowbreak{}h and r=\allowbreak{}+\allowbreak{}/\allowbreak{}-j modulo 2h. Since 2\textasteriskcentered{}\allowbreak{}chi\_\allowbreak{}a(h-1)\textless{}\allowbreak{}chi\_\allowbreak{}a(h), two coordinates below h cannot add to u. Congruence up to sign forces r=\allowbreak{}i. The endpoint i=\allowbreak{}h is the unique zero residue. In particular congruence to +\allowbreak{}/\allowbreak{}-a forces j odd.

\end{enumerate}



These statements are unbounded arithmetic lemmas; finite Pell tests do not replace their proofs. Polynomial dependence of chi\_\allowbreak{}b(j),psi\_\allowbreak{}b(j) on b permits base transfer modulo any integer. At b=\allowbreak{}1 one has chi\_\allowbreak{}1(j)=\allowbreak{}1, psi\_\allowbreak{}1(j)=\allowbreak{}j.



\subsection*{3. A common square frame and common nonzero ordinate}



Let Aset be the distinct actual base, exponent, and output atoms of the eighteen original power requests, using their exact affine restorations in the contracted source. Replace the private frame by



\begin{verbatim}
W=2*(2+sum(a^2 for a in Aset)).
\end{verbatim}

Consequently W\textgreater{}\allowbreak{}=\allowbreak{}4, |a|\textless{}\allowbreak{}W, a\textasciicircum{}\allowbreak{}2\textless{}\allowbreak{}W, and 2|a|\textless{}\allowbreak{}W for each atom. In particular the requested doubled exponents 2k and squared outputs m\textasciicircum{}\allowbreak{}2 are strictly below W. No sign of a proposed power output is needed for this frame. Its exact affine atom list and defining row are in the transformation certificate.



Retain A\textgreater{}\allowbreak{}=\allowbreak{}0 and Zminus1\textgreater{}\allowbreak{}=\allowbreak{}0 and the old frame norm



\begin{verbatim}
A^2-W^3(W+2)*(Zminus1+1)^2=1.
\end{verbatim}

Set Z=\allowbreak{}Zminus1+\allowbreak{}1\textgreater{}\allowbreak{}=\allowbreak{}1. Pell classification at center W+\allowbreak{}1 identifies (A,WZ)=\allowbreak{}(chi\_\allowbreak{}(W+\allowbreak{}1)(h),psi\_\allowbreak{}(W+\allowbreak{}1)(h)) with h\textgreater{}\allowbreak{}=\allowbreak{}1. Reduction modulo W gives psi\_\allowbreak{}(W+\allowbreak{}1)(h)=\allowbreak{}h modulo W, so W divides h and h\textgreater{}\allowbreak{}=\allowbreak{}W. Therefore



\begin{verbatim}
A >= (W+1)^W > W^W,  A>W^2.
\end{verbatim}

For every positive actual power base 1\textless{}\allowbreak{}=\allowbreak{}n\textless{}\allowbreak{}W, the existing definition



\begin{verbatim}
T_n=2An-n^2-1
\end{verbatim}

gives T\_\allowbreak{}n\textgreater{}\allowbreak{}=\allowbreak{}2A-2\textgreater{}\allowbreak{}A\textgreater{}\allowbreak{}W\textasciicircum{}\allowbreak{}W. This follows because the quadratic increases on 1\textless{}\allowbreak{}=\allowbreak{}n\textless{}\allowbreak{}W\textless{}\allowbreak{}A and its minimum is at n=\allowbreak{}1.



Keep the common norm U\textasciicircum{}\allowbreak{}2-(A\textasciicircum{}\allowbreak{}2-1)Vaux\textasciicircum{}\allowbreak{}2=\allowbreak{}1, the original product identity



\begin{verbatim}
Vaux=(product over fifteen seeds Y_i^2)*(v+1),
\end{verbatim}

and impose Vaux\textgreater{}\allowbreak{}=\allowbreak{}1 directly. Delete the old sign predicate v\textgreater{}\allowbreak{}=\allowbreak{}0. The product and positivity imply all Y\_\allowbreak{}i are nonzero, their squares divide Vaux, and v+\allowbreak{}1\textgreater{}\allowbreak{}0, hence v\textgreater{}\allowbreak{}=\allowbreak{}0. Thus this change supplies its own omitted sign before using the index theorem. The common norm implies |U|\textgreater{}\allowbreak{}=\allowbreak{}A\textgreater{}\allowbreak{}=\allowbreak{}2.



Retain the constructor Baux=\allowbreak{}1+\allowbreak{}4Vaux\textasteriskcentered{}\allowbreak{}(q4+\allowbreak{}1) with q4\textgreater{}\allowbreak{}=\allowbreak{}0 and Baux=\allowbreak{}A modulo U. Thus Baux\textgreater{}\allowbreak{}=\allowbreak{}5 is odd, and Baux=\allowbreak{}1 modulo 4|Y\_\allowbreak{}i| for every seed. No auxiliary comparison coordinate or congruence quotient needs a sign assumption.



\subsection*{4. Even seed indices recover their own missing size bounds}



For a selected seed, k\textgreater{}\allowbreak{}=\allowbreak{}1 is its positive original exponent and e=\allowbreak{}2k\textless{}\allowbreak{}W. Keep its three norms and abscissa congruence, but change its ordinate congruence to



\begin{verbatim}
X^2-(A^2-1)Y^2=1,
S^2-(Baux^2-1)Raux^2=1,
S=X (mod U),  Raux=2k (mod 4Y).
\end{verbatim}

The private old row Y=\allowbreak{}k+\allowbreak{}slack may be removed with its unconstrained slack. It has no other occurrence after its sign guard is removed. Its unique inverse is slack=\allowbreak{}Y-k; this exact pruning is separately certified.



We already know Y!=\allowbreak{}0. Suppose |Y|=\allowbreak{}1. Then |X|=\allowbreak{}A. Classify the comparison norm by S=\allowbreak{}+\allowbreak{}/\allowbreak{}-chi\_\allowbreak{}Baux(j), Raux=\allowbreak{}+\allowbreak{}/\allowbreak{}-psi\_\allowbreak{}Baux(j). Transfer Baux=\allowbreak{}A modulo U and S=\allowbreak{}X modulo U. The abscissa residue fact implies that j is odd. Since Baux is odd, psi\_\allowbreak{}Baux(j)=\allowbreak{}j modulo 2; hence Raux is odd. But Raux=\allowbreak{}2k modulo 4Y makes it even. Contradiction. Therefore |Y|\textgreater{}\allowbreak{}=\allowbreak{}2. Its main Pell index i is at least two, and



\begin{verbatim}
|Y|=psi_A(i)>=psi_A(2)=2A>2k.
\end{verbatim}

Now classify the common norm as |U|=\allowbreak{}chi\_\allowbreak{}A(h), Vaux=\allowbreak{}psi\_\allowbreak{}A(h), h\textgreater{}\allowbreak{}=\allowbreak{}1. Squared divisibility gives i\textless{}\allowbreak{}=\allowbreak{}h and |Y| divides h. The S congruence yields j=\allowbreak{}+\allowbreak{}/\allowbreak{}-i modulo 2h. Since Baux=\allowbreak{}1 modulo 4|Y|, the Raux congruence yields j=\allowbreak{}+\allowbreak{}/\allowbreak{}-2k modulo 4|Y|. Therefore i=\allowbreak{}+\allowbreak{}/\allowbreak{}-2k modulo 2|Y|. But 1\textless{}\allowbreak{}=\allowbreak{}i\textless{}\allowbreak{}=\allowbreak{}|Y| and 1\textless{}\allowbreak{}=\allowbreak{}2k\textless{}\allowbreak{}|Y|; the minus congruence is impossible and the plus congruence forces i=\allowbreak{}2k. Thus



\begin{verbatim}
|X|=chi_A(2k), |Y|=psi_A(2k).
\end{verbatim}

For the two undoubled t and K seeds, their retained guards Y\textgreater{}\allowbreak{}=\allowbreak{}k\textgreater{}\allowbreak{}=\allowbreak{}1 give the usual same short-interval proof, with positive Y and |X|=\allowbreak{}chi\_\allowbreak{}A(k). All exponent positivity is syntactic before power semantics: t\textgreater{}\allowbreak{}=\allowbreak{}16, K=\allowbreak{}2R\textgreater{}\allowbreak{}0, t+\allowbreak{}F\textasciicircum{}\allowbreak{}2\textgreater{}\allowbreak{}0, x+\allowbreak{}1\textgreater{}\allowbreak{}0, y+\allowbreak{}1\textgreater{}\allowbreak{}0, and every encoded representative plus one is positive. The derived label J, when present, has a monotone constructor with positive constant term. Its source certificate supplies the compositional lower bound.



\subsection*{5. Squared outputs force positive Pell coordinates}



For a doubled seed and a supporting base n\textgreater{}\allowbreak{}=\allowbreak{}2, replace its power congruence by



\begin{verbatim}
X+(n-A)Y=m^2+T_n*h.
\end{verbatim}

Set e=\allowbreak{}2k, N0=\allowbreak{}n\textasciicircum{}\allowbreak{}e, M=\allowbreak{}m\textasciicircum{}\allowbreak{}2, x=\allowbreak{}chi\_\allowbreak{}A(e), y=\allowbreak{}psi\_\allowbreak{}A(e). The standard recurrence and norm give



\begin{verbatim}
H=x+(n-A)y = N0 (mod T_n),
J=x-(n-A)y,  HJ=1+T_n*y^2.
\end{verbatim}

The four independent signs of X,Y imply respectively that T\_\allowbreak{}n divides



\begin{verbatim}
M-N0,  M+N0,  M*N0-1,  M*N0+1.
\end{verbatim}

The frame has n,e,M\textless{}\allowbreak{}W. Hence N0\textless{}\allowbreak{}W\textasciicircum{}\allowbreak{}(W-1), M\textasteriskcentered{}\allowbreak{}N0\textless{}\allowbreak{}W\textasciicircum{}\allowbreak{}W\textless{}\allowbreak{}A\textless{}\allowbreak{}T\_\allowbreak{}n, and M+\allowbreak{}N0\textless{}\allowbreak{}W+\allowbreak{}W\textasciicircum{}\allowbreak{}(W-1)\textless{}\allowbreak{}W\textasciicircum{}\allowbreak{}W\textless{}\allowbreak{}T\_\allowbreak{}n. The second and fourth divisibilities are impossible by strict positivity. The third is impossible because N0\textgreater{}\allowbreak{}=\allowbreak{}4 and M is a nonnegative integer, so M\textasteriskcentered{}\allowbreak{}N0 cannot equal one. Only X=\allowbreak{}x,Y=\allowbreak{}y remains, and the first congruence gives M=\allowbreak{}N0 exactly. Thus m=\allowbreak{}+\allowbreak{}/\allowbreak{}-n\textasciicircum{}\allowbreak{}k and the actual main Pell coordinates are positive, even before the sign of m is resolved.



For q-power outputs already known positive from their remainder bounds, this immediately gives the intended power. For terminal successor-base outputs, Section 1.1 gives the positive sign once their front division yields m=\allowbreak{}1 modulo q. No factorial or binomial value is assumed to obtain this orientation.



\subsection*{6. The noncircular soundness order}



The exact sign-inference certificate uses only integer polynomial identities, positive monomials, and identified affine squares. Its omitted paired predicates are restored first by the weighted interval lemma. The needed q\textgreater{}\allowbreak{}0 is independent: the unchanged factorial remainder/\allowbreak{}slack predicates give



\begin{verbatim}
B=fact_remainder+fact_slack+1>=1,
q=B+choose_oslack+1>=2.
\end{verbatim}

The retained anchor for the binomial remainder gives



\begin{verbatim}
L=q*u+choose_rslack+2>=2.
\end{verbatim}

Every other interval anchor is positive in the same way. No power is interpreted in these statements. After all seven pair restorations, the emitted algebraic closure has every required ordinary front sign. Its selected Y-slacks, comparison coordinates, ordinary division quotients, and power outputs need no such sign.



\begin{enumerate}

\item The retained t-seed guard identifies Y=\allowbreak{}psi\_\allowbreak{}A(t)\textgreater{}\allowbreak{}0 and X=\allowbreak{}+\allowbreak{}/\allowbreak{}-chi\_\allowbreak{}A(t). Its request at base q has output L\textgreater{}\allowbreak{}=\allowbreak{}2. The signed power residue excludes the negative X: multiplying by q\textasciicircum{}\allowbreak{}t would force L\textasteriskcentered{}\allowbreak{}q\textasciicircum{}\allowbreak{}t=\allowbreak{}-1 under the strict frame bound, impossible. Hence L=\allowbreak{}q\textasciicircum{}\allowbreak{}t and this t pair is positively oriented.

\item The other t request now gives V0=\allowbreak{}(F\textasciicircum{}\allowbreak{}2+\allowbreak{}1)\textasciicircum{}\allowbreak{}t\textgreater{}\allowbreak{}0, even if the base was initially one. The literal ceiling F\textasteriskcentered{}\allowbreak{}B=\allowbreak{}V0+\allowbreak{}fact\_\allowbreak{}slack and B\textgreater{}\allowbreak{}=\allowbreak{}1 imply F\textgreater{}\allowbreak{}=\allowbreak{}1. Moreover V0=\allowbreak{}(F\textasciicircum{}\allowbreak{}2+\allowbreak{}1)\textasciicircum{}\allowbreak{}t\textgreater{}\allowbreak{}=\allowbreak{}F\textasciicircum{}\allowbreak{}2+\allowbreak{}1, so F\textasteriskcentered{}\allowbreak{}B\textgreater{}\allowbreak{}=\allowbreak{}F\textasciicircum{}\allowbreak{}2+\allowbreak{}1 and B\textgreater{}\allowbreak{}F\textgreater{}\allowbreak{}=\allowbreak{}1. Thus B\textgreater{}\allowbreak{}=\allowbreak{}2 and q\textgreater{}\allowbreak{}=\allowbreak{}3. This precedes interpreting B as a binomial.

\item The retained K-seed guard and request N=\allowbreak{}B\textasciicircum{}\allowbreak{}K at base B\textgreater{}\allowbreak{}=\allowbreak{}2 orient that pair by the ordinary signed-output small-residue argument. Its other request is then q=\allowbreak{}N\textasciicircum{}\allowbreak{}K at positive base N\textgreater{}\allowbreak{}=\allowbreak{}2. Both equalities are now genuine.

\item The derived t+\allowbreak{}1 pair is still the original one-step addition from the positively oriented t pair. Its q\textasciicircum{}\allowbreak{}(t+\allowbreak{}1) output is correct. The two selected q-power seeds use Section 5 and their positive interval outputs to recover q\textasciicircum{}\allowbreak{}(y+\allowbreak{}1) and q\textasciicircum{}\allowbreak{}(x+\allowbreak{}1). All front divisors L are therefore positive powers of q of exponent at least one, and hence divisible by q.

\item Apply Sections 4-5 to the remaining eleven selected seeds. Their outputs have |m|=\allowbreak{}(q+\allowbreak{}1)\textasciicircum{}\allowbreak{}k. Every front digit/\allowbreak{}division equation gives m=\allowbreak{}1 modulo q through the explicit substituted remainder. Since q\textgreater{}\allowbreak{}=\allowbreak{}3, the negative sign is impossible. All eighteen original power equalities have been recovered.

\item Restore each original front remainder r=\allowbreak{}1+\allowbreak{}q\textasteriskcentered{}\allowbreak{}u\textgreater{}\allowbreak{}=\allowbreak{}1. Its exact old bound is r+\allowbreak{}s+\allowbreak{}1=\allowbreak{}L, so it is nonnegative and bounded. The original high-digit and range quotients, though initially signed, are nonnegative whenever required: m=\allowbreak{}L(C+\allowbreak{}q\textasteriskcentered{}\allowbreak{}h)+\allowbreak{}r, 0\textless{}\allowbreak{}=\allowbreak{}C\textless{}\allowbreak{}q, 0\textless{}\allowbreak{}=\allowbreak{}r\textless{}\allowbreak{}L and m\textgreater{}\allowbreak{}0 exclude h\textless{}\allowbreak{}=\allowbreak{}-1; similarly m=\allowbreak{}N\textasteriskcentered{}\allowbreak{}L\textasteriskcentered{}\allowbreak{}h+\allowbreak{}r with N\textgreater{}\allowbreak{}=\allowbreak{}1 excludes h\textless{}\allowbreak{}=\allowbreak{}-1. All remaining front meanings follow in the source's factorial, exact-binomial, prime-power lifting, CRT-range, and direct-calculus order.

\end{enumerate}



For reference, the ordinary signed-output exclusion used in steps 1 and 3 needs no sign of the candidate m. With positive Y=\allowbreak{}psi\_\allowbreak{}A(k), negative X would give m\textasteriskcentered{}\allowbreak{}n\textasciicircum{}\allowbreak{}k=\allowbreak{}-1 after the small-residue argument. This is impossible for n\textgreater{}\allowbreak{}=\allowbreak{}2. The frame bounds |m|,n,k\textless{}\allowbreak{}W ensure |m\textasteriskcentered{}\allowbreak{}n\textasciicircum{}\allowbreak{}k+\allowbreak{}1|\textless{}\allowbreak{}T\_\allowbreak{}n. Positive X then gives m=\allowbreak{}n\textasciicircum{}\allowbreak{}k. The exceptional base n=\allowbreak{}1 is used only after the pair is already oriented.



This recovers a genuine semantic front tuple. If the source theorem requires its old private Pell witnesses rather than just the semantic interface, rebuild fresh ones at its original indices and frame by Section 7. No pointwise claim that doubled Pell coordinates solve an undoubled system is made.



\subsection*{7. Complete witness construction in the reverse direction}



Begin with a genuine semantic front tuple supplied by the inherited natural source, not with a hypothetical standard-model consistency assumption. Keep its logical/\allowbreak{}arithmetic coordinates, and perform the following finite choices.



\begin{enumerate}

\item For each terminal successor-base division, take its genuine old remainder r and set u=\allowbreak{}(r-1)/\allowbreak{}q. As proved in Section 1.1, this is a nonnegative integer. Keep its old slack s. The new front equations and weighted sign predicates hold. Choose the frame W from the original base/\allowbreak{}exponent/\allowbreak{}output atoms.

\item Put A=\allowbreak{}chi\_\allowbreak{}(W+\allowbreak{}1)(W), Z=\allowbreak{}psi\_\allowbreak{}(W+\allowbreak{}1)(W)/\allowbreak{}W. The recurrence modulo W proves that Z is a positive integer. Set Zminus1=\allowbreak{}Z-1. This supplies the common frame norm and all positive moduli T\_\allowbreak{}n.

\item For each of the two undoubled seeds choose (X,Y)=\allowbreak{}(chi\_\allowbreak{}A(k),psi\_\allowbreak{}A(k)); for each of thirteen selected seeds choose (chi\_\allowbreak{}A(2k),psi\_\allowbreak{}A(2k)). Derive the t+\allowbreak{}1 pair by the unchanged one-step recurrence. Let P be the product of all fifteen seed ordinates squared.

\item In the finite ring (Z/\allowbreak{}(4P)Z)[s]/\allowbreak{}(s\textasciicircum{}\allowbreak{}2-(A\textasciicircum{}\allowbreak{}2-1)), A+\allowbreak{}s is invertible, with inverse A-s. A positive power is one. Its integer Pell coordinates give U=\allowbreak{}1 modulo 4P, Vaux=\allowbreak{}0 modulo 4P, U\textgreater{}\allowbreak{}0, Vaux\textgreater{}\allowbreak{}0. Thus P divides Vaux, and v=\allowbreak{}Vaux/\allowbreak{}P-1\textgreater{}\allowbreak{}=\allowbreak{}0. The Pell norm gives gcd(U,Vaux)=\allowbreak{}1; U is odd, so gcd(U,4Vaux)=\allowbreak{}1.

\item CRT supplies Baux=\allowbreak{}A modulo U, Baux=\allowbreak{}1 modulo 4Vaux. Increase it along that progression until Baux\textgreater{}\allowbreak{}A and Baux\textgreater{}\allowbreak{}1+\allowbreak{}4Vaux. Its two exact quotient witnesses and q4 are then nonnegative. For each actual requested index e, set S=\allowbreak{}chi\_\allowbreak{}Baux(e), Raux=\allowbreak{}psi\_\allowbreak{}Baux(e). Polynomial base transfer supplies S=\allowbreak{}X modulo U and Raux=\allowbreak{}e modulo 4Y. Their quotient witnesses are integers.

\item Every genuine output makes its unchanged or squared power congruence exact. For e=\allowbreak{}k or 2k the identity

\end{enumerate}



\begin{verbatim}
   chi_A(e)+(n-A)psi_A(e)-n^e
     =T_n*sum_{j=0}^{e-1} n^(e-1-j)*psi_A(j)
\end{verbatim}

supplies the integer quotient directly. Fill every ordinary defining gate by its displayed polynomial. The retained t/\allowbreak{}K index slacks are nonnegative. Fill all retained sign triples using the three-square representation.



The choice order is front data, W, A and the main pairs, the common U,Vaux, Baux, then comparison pairs and quotients. No later-chosen modulus is already constrained by a residue depending on its future value. The common finite-ring modulus contains all fifteen ordinate squares, so these choices are simultaneous. This proves completeness. Using undoubled indices and the old frame gives the fresh old certificate needed in the soundness map.



\subsection*{8. Exact algebra after the semantic change}



The guard expansion and private-slack pruning leave a quadratic integer system. The old Horn product is then factored exactly:



\begin{verbatim}
(product L_i^2)*(Bguard^2+G^2)
  =(product L_i)^2*(Bguard^2+G^2).
\end{verbatim}

The signed product is squared before being added to the two independent beta squares; no cancellation argument is used. The guarded factor ((c-a)\textasciicircum{}\allowbreak{}2\textasteriskcentered{}\allowbreak{}(c-b)\textasciicircum{}\allowbreak{}2-(d+\allowbreak{}1)\textasciicircum{}\allowbreak{}2)\textasciicircum{}\allowbreak{}2 is likewise computed by first forming the signed product (c-a)(c-b). The existing round14 factorization certificate records all changed definitions and verifies that no altered intermediate escapes its subnetwork. Its proof applies unchanged to the new front source.



Next, monic integral projection deletes only actual rows z-Q=\allowbreak{}0 (up to sign), with integral Q of degree at most two. Every substitution is made everywhere and checked not to increase any residual degree above two. A quadratic Q is substituted only where z has no nonlinear occurrence outside its defining row. Reverse substitution restores every deleted coordinate; unused coordinates may be restored as zero. The separately pruned Y-slacks restore as Y-k.



Finally the final variables undergo integral translations. In particular old t0=\allowbreak{}new\_\allowbreak{}t-16, not new\_\allowbreak{}t-1, absorbs the source's new lower-bound offset. The additional recorded finite centering reduces coefficient height and support. Every final shift is explicitly listed; old=\allowbreak{}new+\allowbreak{}shift and new=\allowbreak{}old-shift are inverse maps on all integer tuples. These translations, unlike the remainder restriction in Section 1, are genuine integral bijections. A deterministic finite coordinate-descent search found them; no optimality is claimed.



The final polynomial is the sum of the squares of the resulting residuals. An integer tuple makes that sum zero iff every residual is zero. The whole collected expansion is regenerated exactly. The highest occurring degree is four, and every retained variable occurs in the supplied polynomial.



\subsection*{9. Reproduction, dependency boundary, and credit}



Run the following from the round15 package root:



\begin{verbatim}
python direct/build_residue.py architecture/unified_derived_shell/raw direct/rebuilt_source --reparam
python direct/center_integer.py direct/rebuilt_source direct/rebuilt_integer442
\end{verbatim}

Only Python's standard library is required. Inputs are literal source JSON; no upstream historical builder, network fetch, theorem prover, or solution search is executed. The compiler records source hashes, all selected-index changes, frame atoms, complete remainder substitutions, guard predicates, weighted-pair anchors, exact private-slack inverses, Horn factorization, monic projection, and final translations. The polynomial coefficients are reconstructible from the printed residual specification and from these files.



The even-index parity argument and shared Pell completeness are credited to the incoming Carriers and Even Indices package. Ordinary signed-output recovery, square budgets, Horn factorization and exact sign closure are inherited from round14 and related signed-power work. Shared remainder predicates are credited to incoming Structural Compression; the new weighted version accompanies the new quotient remainder. The unified factorial/\allowbreak{}product architecture has its own attribution and proof in the architecture packet. The new contribution here is their selective composition with residue-one orientation, which avoids a new nonnegativity triple for each signed terminal power, and its complete smaller integer quartics. Novelty and comparisons are relative to inspected sources.



Finite checks certify literal algebra and source binding. The unbounded Pell, integer sign, remainder, factorial, CRT, and refutation arguments prove the semantic equivalence. The named inherited formal calibration is retained as a dependency; its full transitive theorem-prover closure was not rebuilt here.
